\begin{frame}{Un servicio web: una función que otro sistema llama por la red}
\begin{columns}[c]
\begin{column}{0.46\textwidth}
\begin{itemize}
\item Expone operaciones sobre HTTP, con un contrato que dice qué se envía y qué se recibe.
\item Al cliente no le importa el lenguaje ni la máquina del servidor.
\item Dos estilos: \alert{SOAP} (mensajes XML) y \alert{REST} (recursos y JSON).
\end{itemize}
\end{column}
\begin{column}{0.5\textwidth}
\centering
\begin{tikzpicture}[caja/.style={draw=verde,rounded corners=3pt,fill=verdesuave,text width=2.3cm,minimum height=0.9cm,font=\scriptsize,align=center}, flecha/.style={-{Stealth[length=2.5mm]},thick,color=tinta}, etiqueta/.style={font=\scriptsize,color=gris}]
\node[caja] (conta) {Software contable\\(cualquier lenguaje)};
\node[caja,right=1.5cm of conta] (sri) {SRI: recepción\\y autorización};
\node[caja,below=0.9cm of conta] (ang) {Angular\\(TypeScript)};
\node[caja,right=1.5cm of ang] (spring) {Spring Boot\\(Java)};
\draw[flecha] (conta) -- node[above,etiqueta]{SOAP / XML} (sri);
\draw[flecha] (ang) -- node[above,etiqueta]{REST / JSON} (spring);
\end{tikzpicture}
\end{column}
\end{columns}
\fuente{w3c_soap12,fielding2000,sri_recepcion}
\end{frame}

\begin{frame}{¿A qué escala se usan?}
\begin{columns}[T]
\begin{column}{0.31\textwidth}
\cifra{192 millones}{de peticiones por segundo a la API de DynamoDB en Prime Day 2026, y más de 59 billones en cuatro días \cite{aws_primeday2026}}
\end{column}
\begin{column}{0.31\textwidth}
\cifra{500}{API REST públicas halladas entre los 4000 sitios más visitados; solo el 0,8\,\% cumple todos los principios REST \cite{neumann2021}}
\end{column}
\begin{column}{0.31\textwidth}
\cifra{SOAP}{en Ecuador: los comprobantes electrónicos se envían a los servicios de recepción y autorización del SRI \cite{sri_recepcion,sri_autorizacion}}
\end{column}
\end{columns}
\vspace{16pt}
\begin{block}{En nuestro entorno}
Nuestros tres proyectos ya se comunican por servicios web: el front en Angular le pide los datos al backend por REST. En la segunda parte cada uno muestra el suyo.
\end{block}
\end{frame}
